\begin{statement}{C.2.UN.84}{Proved}{The highest Hopf order can be made the global maximum}
Let \(f:X\to\Delta\) be a proper degeneration of connected
complex two-tori with smooth total space and full reduced
SNC central fibre. Use its actual primitive canonical
generator and write
\[
 \operatorname{div}(\Omega)=\sum_i k_iV_i,\qquad
 \min_i k_i=0,\qquad R=\max_i k_i.
 \tag{C.TK.1}
\]
Let \(\mathcal H\) be the set of prime components whose
smooth minimal model is Hopf, and suppose that
\[
                  h=\max_{i\in\mathcal H}k_i>0.
 \tag{C.TK.2}
\]
There is a finite overlap-compatible modification record,
over the same disc and isomorphic on its punctured part,
whose endpoint has smooth total space and full reduced
NC fibre, such that its global maximum canonical order
is exactly \(h\). Every original Hopf-birational prime
survives with its original order, and no new such prime
appears. Every component of order \(h\) with Hopf minimal
model can, in the same finite record, be made an embedded
minimal Hopf surface. These maximal Hopfs are the
torsion-normal elliptic leaves of UN82.

At all retained global stages the canonical orders lie
between zero and the original \(R\). A whole-component
removal has positive order \(q\); its possible added
components have order \(q+1\), bounded by the current
maximum, and there are no additions when \(q\) is that
maximum. The removed \(q\) is not required to exceed
\(h\). The termination potential therefore counts every
positive order, including orders at most \(h\).

The record retains its actual centres, normalization
branches, local sigma maps and restored overlaps.
Reversing it recovers the given analytic family. On
an original compact family it gives the same actual
modification over its compact base. Its Euler changes
are the signed normal-cone terms of UN75.
\end{statement}

\begin{proof}
\par\medskip\noindent\textbf{The precise input to the component selection.}
We use the actual selection argument, rather than
applying the conclusion of Kulikov's canonical-flat
theorem to a family that does not satisfy its
Hopf-minimum hypothesis. The relevant distinction
is explicit in \cite[\S7]{Kulikov1981}: the search
for a component whose normalization is ruled with
one boundary section uses only that the current
maximal components are ruled. Section8 removes
the earlier dual-complex hypothesis and supplies
the missing-fibre blowups. Thus a lower-order Hopf
component is permitted during this selection.

Here are the actual maps and the order information
from that argument which we retain. The boundary
analysis at the maximal ruled components in
\cite[\S\S4.8.5--4.8.12,4.16.2--4.16.6]{Kulikov1977}
produces either a plane with one line or a ruled
normalization with one section and vertical boundary.
The two-section cases are converted by the specified
finite type-I/type-II moves along their section
chains or cycles; a rational tree gives the complete
one-triple ruling curve. At a singular sigma chart,
the last locally blown-up branch of largest order
supplies the same one-section configuration. These
are the actual two-ruling, complete-curve and
sigma-atlas maps recorded in UN75 and UN77.

The ruled candidate need not be maximal. If its
order is \(q\), restriction of the common canonical
section to a generic complete ruling fibre gives
the section coefficient \(-2\), hence
\[
       k_{\mathrm{section\ parent}}=q-1.
 \tag{C.TK.3}
\]
For a plane candidate, the line coefficient \(-3\)
gives the parent order \(q-2\). Consequently every
removed component has positive order. When a ruled
candidate has disconnected boundary, the corrected
procedure inserts the missing fibre components.
Their actual curve-blowup Jacobians give order
\(q+1\). If \(q\) is the current maximum, the
effective anticanonical one-section comparison
of \cite[Lemma2.19]{Kulikov1977} makes the boundary
connected and no insertion is made. In every case
\[
 \begin{split}
 &\text{one component of order }q>0\text{ is removed},\\
 &\text{only finitely many components of order }q+1
       \text{ may be added},\\
 &q+1\le R_{\mathrm{current}}\text{ whenever an addition occurs}.
 \end{split}
 \tag{C.TK.4}
\]
This is an assertion about the retained degeneration
stages. It does not assign these bounds to auxiliary
common blowups in a two-ruling diagram.

The local contractions are followed by the exact
restoration of the overlap, using the same composite
map. The proper-graph comparison in (C.KG.2) and
(C.KG.6) therefore gives an actual bimeromorphic
map over the disc, not merely equal central graphs.
If singular sigma charts remain, we use the
largest last locally blown-up branch and its finite
descendant cleanup from UN77. Such branches are
ruled by their actual complete exceptional fibres.
This cleanup uses no hypothesis on the other
Hopf orders and has the same positive-level updates
(C.TK.4). It ends at a smooth-total NC stage with
the labelled normalization branches retained.

\par\medskip\noindent\textbf{Why selection remains available above \(h\).}
The proper composite-graph argument in UN75 proves
that every Hopf-birational prime survives every
recorded removal. A general complete rational
exceptional fibre cannot cover a point-blown-up
Hopf surface; a compact projective point exceptional
cannot map bimeromorphically onto it. The actual
canonical direct-image isomorphism (C.KG.6) preserves
its canonical order. Every added prime is ruled or
projective over the specified centre, so it is not
Hopf-birational. It follows that the set of Hopf
orders, with its original labels, is unchanged.
In particular its highest value remains \(h\).

Consider a current smooth-total NC stage with
maximum \(R'>h\). On the smooth normalization
\(T\) of a maximal component, actual adjunction is
\[
       -K_T=\sum_C (R'-k_{\mathrm{other}(C)}+1)C.
 \tag{C.TK.5}
\]
Both preimages of a self-double curve are counted;
their coefficients are one. The divisor is nonzero:
the connected support contains a component of Hopf
order \(h<R'\), so no maximal component is an
isolated component of the whole support. Every
coefficient in (C.TK.5) is positive. A positive
pluricanonical section would be a holomorphic
function on the compact connected surface vanishing
on this effective divisor, which is impossible.
Thus \(\kappa(T)=-\infty\).

The actual component theorem UN19 makes a
nonruled component of this Kodaira dimension
Hopf-birational. No such component has order
\(R'>h\). Hence every current maximal
normalization is ruled or rational, including
the possible plane. This is exactly the input
to the selection proved above. Controlled
singular sigma charts are first handled by
their specified last-local ruled branches;
they do not require an anticanonical formula
on a singular ambient chart.

\par\medskip\noindent\textbf{The stopped finite-level induction.}
For each \(1\le j\le R\), let \(c_j\) be the
number of retained global components of order
\(j\), and order the tuples
\[
                     (c_1,c_2,\ldots,c_R)
 \tag{C.TK.6}
\]
lexicographically, with the smaller index first.
Each whole-component block (C.TK.4) strictly
decreases this tuple: it deletes a card of
order \(q\), leaves all smaller orders unchanged,
and only increases the next order. Type-I and
type-II preparatory moves preserve the card
counts. A fixed finite product of copies of
\(\mathbb N\) in this order has no infinite
strict descending sequence. The finite cleanup
blocks obey the same comparison.

We continue only while the maximum exceeds
\(h\), and finish the finite cleanup belonging
to the last selected block. If the maximum
still exceeded \(h\), (C.TK.5) would supply
another selection. The decreasing tuple
therefore forces a finite stop with maximum
at most \(h\). The retained Hopf prime of
order \(h\) forces equality. No order-zero
prime has been removed, by (C.TK.3) and its
plane counterpart, so the transported canonical
generator remains primitive with minimum zero.

Counting only the components above \(h\)
would not prove this induction. For example,
removing an order-one ruled candidate and
adding order-two components changes none of
those upper counts when \(h\ge2\). Such a
candidate can be the ruled neighbour produced
by the one-triple step. The full tuple
(C.TK.6) is strictly decreased by the exact
update and is the potential actually used.

\par\medskip\noindent\textbf{The actual maximal Hopf components at the stopping stage.}
The original central components were embedded.
The record has introduced no elliptic self-double
curve on a surviving Hopf component. Type-I and
type-II changes create incidences on their
rational exceptional curves. A removed ruled
component has one horizontal boundary section;
its generic contraction cannot identify two
elliptic horizontal branches. Its other changes
are on its rational fibre curves, and its
restored overlaps retain the original elliptic
maps. Point contractions likewise only change
the indicated rational boundary configuration.

At the stopping stage a Hopf normalization
of order \(h\) has the effective divisor
(C.TK.5), now with \(R'=h\). The actual
point-map induction (C.HM.3)--(C.HM.4)
shows that a rational coefficient-one boundary
is a leaf, and that two coefficient-one
boundaries never meet. A rational self-double
curve would have two coefficient-one preimages,
each with negative self-square and exactly
one boundary intersection. Their actual normal
product would consequently give
\[
                        C_1^2+C_2^2=-1,
 \tag{C.TK.7}
\]
contradicting the two negative integer squares.
A triple point with three branches of this same
component would already give two intersecting
coefficient-one curves. Thus these maximal
Hopfs are embedded at the stopped NC stage.

All the prerequisites of the last-point and
finite-cleanup proof of UN77 now hold on this
actual recorded stage. That proof uses its
smooth normalizations, the labelled NC branches
elsewhere and the embedded maximal Hopfs; it
does not require the other components to be
embedded. Apply those exact moves to remove
their point defects. The retained global orders
stay at most \(h\), every Hopf prime and order
is preserved, and the endpoint maximal Hopfs
are smooth minimal surfaces. UN82 then supplies
their single elliptic boundaries and exact
torsion-normal line maps.

Finally the local proper graphs glue with the
recorded overlap isomorphisms, exactly as in
UN75. They give the punctured-family isomorphism
and the inverse record on its actual endpoint
neighbourhood. Annular gluing to the unchanged
outside performs the same operation on an
original compact family. The signed normal-cone
Euler formula (C.KG.3)--(C.KG.4) applies to
every retained arrow. No chosen smoothing of
the endpoint central space replaces that
actual neighbourhood.
\end{proof}


% C.2.UN.1 from part_C2_unrestricted.tex

